% titel: Blatt 2, A3: Das mathematische Pendel
% skript: Kap. Schwingungen, mathematisches Pendel

% Bewegungsgleichung und Schwingungsdauer des Fadenpendels.

Das mathematische (Faden-)Pendel schwingt fuer kleine Auslenkungen
harmonisch; die Schwingungsdauer haengt nur von Laenge und Ortsfaktor ab,
nicht von der Masse.

\begin{werkzeugkasten}[Werkzeuge]
\textbf{Ruecktreibende Kraft.} Fuer eine um $\varphi$ ausgelenkte
Punktmasse am Faden der Laenge $\ell$ wirkt die Tangentialkomponente der
Gewichtskraft; Newton in Tangentialrichtung ($s=\ell\varphi$):
\begin{equation}
  m\ell\ddot\varphi=-mg\sin\varphi
  \;\Rightarrow\;\;
  \ddot\varphi+\frac{g}{\ell}\sin\varphi=0\;.
\end{equation}

\textbf{Kleine Auslenkungen} $\sin\varphi\approx\varphi$:
\begin{equation}
  \ddot\varphi+\frac{g}{\ell}\,\varphi=0,
  \qquad\omega=\sqrt{\frac{g}{\ell}},
  \qquad T=2\pi\sqrt{\frac{\ell}{g}}\;.
\end{equation}

\textbf{Ansatz} $e^{i\omega t}$ (oder $\cos$, $\sin$): einsetzen ergibt
eine Bedingung fuer $\omega$. Damit laesst sich aus gemessener Periode die
Fadenlaenge $\ell$ zurueckrechnen.
\end{werkzeugkasten}

\begin{werkzeugkasten}[Vorgehen]
\begin{itemize}
  \item Zuerst die Bewegungsgleichung in Winkelkoordination aufstellen.
  \item Die Naherung $\sin\varphi\approx\varphi$ anwenden und $\omega$
        ablesen.
  \item Aus $T$ die Laenge $\ell=g\,T^2/(4\pi^2)$ bestimmen; zur
        Deckenhoehe addiert sich der Abstand der Lampe ueber dem Boden.
\end{itemize}
\end{werkzeugkasten}

\begin{achtung}[Typische Fallen]
\begin{itemize}
  \item $T$ aus mehreren Perioden durch Division bekommen; die angegebene
        Zeit ist die fuer 5 Schwingungsperioden.
  \item $\omega$ ist die \emph{Kreisfrequenz} ($\omega=2\pi/T$), nicht $f$.
  \item Die Deckenhoehe ist $\ell$ \emph{plus} die H\"ohe der Lampe ueber dem
        Boden, nicht nur $\ell$.
\end{itemize}
\end{achtung}
